\subsection{INVERSE AND DIRECT LIMITS OF ALGEBRAS}

Let I be a preordered set and \((\mathbf{A}_{i},\phi_{ij})\) an inverse system of commutative rings with I as indexing set. Let \((\mathbf{E}_{i},f_{ij})\) be an inverse system of \(A_{i}\) -modules with I as indexing set (II, §6, no. 1) and suppose further that each \(E_{i}\) has an \(A_{i}\) -algebra structure and that, for \(i\leqslant j,f_{ij}\) is an \(A_{j}\) -homomorphism of algebras (relative to \(\phi_{ij}\) ) (no. 5). Let \(A=\varprojlim A_{i}\) and \(E=\varprojlim E_{i}\) , which has an A-module structure, the inverse limit of the structure of the \(\mathbf{A}_i\) -modules \(\mathbf{E}_i\) (II, §6, no. 1); it is immediately verified that the law of composition on \(\mathbf{E}\) , considered as the inverse limit of the \(\mathbf{E}_i\) considered as magmas under multiplication (I, §10, no. 1), with the A-module structure on \(\mathbf{E}\) , defines on \(\mathbf{E}\) an A-algebra structure; \((\mathbf{E}_i, f_{ij})\) is called an inverse system of \(\mathbf{A}_i\) -algebras and the A-algebra \(\mathbf{E}\) is called its inverse limit. If \(f_i: \mathbf{E} \to \mathbf{E}_i\) , \(\phi_i: \mathbf{A} \to \mathbf{A}_i\) are the canonical mappings, \(f_i\) is an A-homomorphism of algebras (relative to \(\phi_i\) ). If the \(\mathbf{E}_i\) are associative (resp. commutative), so is \(\mathbf{E}\) ; if each \(\mathbf{E}_i\) admits a unit element \(e_i\) and \(f_{ij}(e_j) = e_i\) for \(i \leqslant j\) , \(e = (e_i)\) is a unit element of the algebra \(\mathbf{E}\) .

Let \((\mathbf{E}_{i}^{\prime}, f_{ij}^{\prime})\) be another inverse system of \(A_{i}\) -algebras and for all i let \(u_{i}: E_{i} \to E_{i}^{\prime}\) be an \(A_{i}\) -algebra homomorphism, these mappings forming an inverse system; then \(u = \lim u_{i}\) is an A-algebra homomorphism.

Suppose now that all the \(A_{i}\) are equal to the same commutative ring A and the \(\phi_{ij}\) to \(Id_{A}\) , so that \(E = \lim E_{i}\) is an A-algebra. Let F be an A-algebra and, for all \(i \in I\) , let \(u_{i}: F \to E_{i}\) be an A-algebra homomorphism such that \((u_{i})\) is an inverse system of mappings; then \(u = \lim u_{i}\) is a homomorphism of the algebra F into the algebra E. Conversely, for every A-algebra homomorphism \(v: F \to E\) , the family of \(v_{i} = f_{i} \circ v\) is an inverse system of A-algebra homomorphisms such that \(v = \lim v_{i}\) . As moreover, writing \(\bar{f}_{ij} = \text{Hom}(1_{\text{F}}, f_{ij})\) , \((\text{Hom}_{\text{A-alg.}}(\text{F}, \text{E}_i), \bar{f}_{ij})\) is clearly an inverse system of sets, it is seen that the above remarks can also be expressed by saying that the canonical mapping \(v \mapsto (f_i \circ v)\) is a bijection

\[
l _ {\mathrm{F}}: \operatorname{Hom} _ {\mathrm{A-alg.}} (\mathrm{F}, \varprojlim \mathrm{E} _ {i}) \rightarrow \varprojlim \operatorname{Hom} _ {\mathrm{A-alg.}} (\mathrm{F}, \mathrm{E} _ {i}).
\]

Moreover, for every A-algebra homomorphism \(w: \mathbf{F} \to \mathbf{F}'\) , the

\[
\bar {w} _ {i} = \operatorname{Hom} (w, 1 _ {\mathrm{E} _ {i}}): \operatorname{Hom} _ {\mathrm{A-alg.}} (\mathrm{F} ^ {\prime}, \mathrm{E} _ {i}) \rightarrow \operatorname{Hom} _ {\mathrm{A-alg.}} (\mathrm{F}, \mathrm{E} _ {i})
\]

form an inverse system of mappings and the diagram

\[
\begin{array}{c} \operatorname{Hom} _ {\mathrm{A-alg.}} (\mathrm{F} ^ {\prime}, \lim \mathrm{E} _ {i}) \xrightarrow {l _ {\mathbf {F} ^ {\prime}}} \lim \operatorname{Hom} _ {\mathrm{A-alg.}} (\mathrm{F} ^ {\prime}, \mathrm{E} _ {i}) \\ \operatorname{Hom} (w, 1 _ {\mathbf {E}}) \Biggl \downarrow \qquad \qquad \qquad \qquad \qquad \qquad \qquad \qquad \qquad \qquad \qquad \qquad \qquad \qquad \qquad \qquad \qquad \qquad \qquad \qquad \qquad \qquad \qquad \qquad \qquad \qquad \qquad \qquad \qquad \qquad \qquad \qquad \qquad \qquad \\ \operatorname{Hom} _ {\mathrm{A-alg.}} (\mathrm{F}, \lim \mathrm{E} _ {i}) \xrightarrow {l _ {\mathbf {F}}} \lim \operatorname{Hom} _ {\mathrm{A-alg.}} (\mathrm{F}, \mathrm{E} _ {i}) \end{array}
\]

is commutative.

Suppose now that I is right directed. Consider a direct system of commutative rings \((\mathbf{A}_i, \phi_{ji})\) and a direct system \((\mathbf{E}_i, f_{ji})\) of \(\mathbf{A}_i\) -modules, with I as indexing set; suppose that each \(\mathbf{E}_i\) has an \(\mathbf{A}_i\) -algebra structure and that, for \(i \leqslant j\) , \(f_{ji}\) is an \(\mathbf{A}_i\) -homomorphism of algebras (relative to \(\phi_{ji}\) ) (no. 5). Let \(\mathbf{A} = \varinjlim \mathbf{A}_i\) , \(\mathbf{E} = \varinjlim \mathbf{E}_i\) ; \(\mathbf{E}\) has an A-module structure, the direct limit of the structures of the \(\mathbf{A}_i\) -modules \(\mathbf{E}_i\) (II, §6, no. 2); moreover, the law of composition on E considered as the direct limit of the \(E_{i}\) , considered as magmas under multiplication (I, §10, no. 3), with the A-module structure on E, defines an A-algebra structure on E; \((\mathbf{E}_{i}, f_{jt})\) is called a direct system of \(A_{i}\) -algebras and the A-algebra E is called its direct limit. If \(f_{i}: E_{i} \to E\) , \(\phi_{i}: A_{i} \to A\) are the canonical mappings, \(f_{i}\) is an \(A_{i}\) -homomorphism of algebras (relative to \(\phi_{i}\) ). If the \(E_{i}\) are associative (resp. commutative), so is E; if each \(E_{i}\) admits a unit element \(e_{i}\) and \(f_{jt}(e_{i}) = e_{j}\) for \(i \leqslant j\) , E admits a unit element e such that \(f_{i}(e_{i}) = e\) for all \(i \in I\) .

Let \((\mathbf{E}_i', f_{ji}')\) be another direct system of \(\mathbf{A}_i\) -algebras and for all \(i\) let \(u_i: \mathbf{E}_i \to \mathbf{E}_i'\) be an \(\mathbf{A}_i\) -algebra homomorphism, these mappings forming a direct system; then \(u = \lim u_i\) is an \(\mathbf{A}\) -algebra homomorphism.

Suppose now that all the rings \(\mathbf{A}_i\) are equal to the same ring \(\mathbf{A}\) and the \(\phi_{ji}\) to \(\mathrm{Id}_{\mathbf{A}}\) , so that \(\mathbf{E} = \varinjlim \mathbf{E}_i\) is an A-algebra. Let \(\mathbf{F}\) be an A-algebra and for all \(i\) let \(u_i: \mathbf{E}_i \to \mathbf{F}\) an A-algebra homomorphism such that \((u_i)\) is a direct system of mappings; then \(u = \varinjlim u_i\) is a homomorphism of the algebra \(\mathbf{E}\) into the algebra \(\mathbf{F}\) . Conversely, for every A-algebra homomorphism \(v: \mathbf{E} \to \mathbf{F}\) , the family of \(v_i = v \circ f_i\) is a direct system of A-algebra homomorphisms such that \(v = \varinjlim v_i\) . As moreover, writing \(\bar{f}_{ij} = \operatorname{Hom}(f_{ij}, 1_F)\) , \((\operatorname{Hom}_{\mathbf{A-alg}}(\mathbf{E}_i, \mathbf{F}), \bar{f}_{ij})\) is clearly an inverse system of sets, it is seen that the above remarks can also be expressed by saying that the canonical mapping \(v \mapsto (v \circ f_i)\) is a bijection

\[
d _ {\mathrm{F}}: \operatorname{Hom} _ {\mathrm{A-alg.}} (\varinjlim \mathrm{E} _ {i}, \mathrm{F}) \rightarrow \varprojlim \operatorname{Hom} _ {\mathrm{A-alg.}} (\mathrm{E} _ {i}, \mathrm{F}).
\]

Further, for every A-algebra homomorphism \(w: F \rightarrow F'\) , the

\[
\bar {w} _ {i} = \operatorname{Hom} (1 _ {\mathrm{E} _ {i}}, w): \operatorname{Hom} _ {\mathrm{A-alg.}} (\mathrm{E} _ {i}, \mathrm{F}) \rightarrow \operatorname{Hom} _ {\mathrm{A-alg.}} (\mathrm{E} _ {i}, \mathrm{F} ^ {\prime})
\]

form an inverse system of mappings and the diagram

\[
\begin{array}{l} \operatorname{Hom} _ {\mathrm{A-alg.}} (\varinjlim \mathrm{E} _ {i}, \mathrm{F}) \xrightarrow {d _ {\mathrm{F}}} \varprojlim \operatorname{Hom} _ {\mathrm{A-alg.}} (\mathrm{E} _ {i}, \mathrm{F}) \\ \operatorname{Hom} (1 _ {\mathrm{E}}, w) \Biggl \downarrow \qquad \qquad \qquad \qquad \qquad \qquad \qquad \qquad \qquad \qquad \qquad \qquad \qquad \qquad \qquad \qquad \qquad \qquad \qquad \qquad \varprojlim \varpi_ {i} \\ \operatorname{Hom} _ {\mathrm{A-alg.}} (\varinjlim \mathrm{E} _ {i}, \mathrm{F} ^ {\prime}) \xrightarrow {d _ {\mathrm{F} ^ {\prime}}} \varprojlim \operatorname{Hom} _ {\mathrm{A-alg.}} (\mathrm{E} _ {i}, \mathrm{F} ^ {\prime}) \end{array}
\]

is commutative.

